\documentclass[10pt,a4paper]{amsart}
\usepackage[margin=19mm]{geometry}
\usepackage{amsmath,amssymb,mathtools}
\usepackage[scr=rsfs,cal=euler]{mathalfa}
\usepackage{xfrac}
\usepackage[unicode,hidelinks,bookmarks=false]{hyperref}
\newcommand{\ie}{i.e.\ }
\newcommand{\cf}{cf.\ }
\newcommand{\resp}{respectively\ }
\newcommand{\Subsection}{\subsection}
\providecommand{\nobreakdash}{\nobreak\hbox{-}}
\setlength{\emergencystretch}{2em}
\multlinegap0pt
\usepackage{amssymb}
\usepackage{enumerate}
\usepackage{mathtools}
\usepackage[shortlabels]{enumitem}
\usepackage{xcolor}
\newtheorem{theorem}{Theorem}
\newtheorem{lemma}{Lemma}
\newcommand{\R}{\mathbb{R}}
\newcommand{\delomega}{\partial\Omega}
\renewcommand{\leq}{\leqslant}
\title{A fractional Michael-Simon Sobolev inequality\\ on convex hypersurfaces}
\author{Gyula Csat\'o, Prosenjit Roy}
\date{Source: arXiv:2407.12439v3 (CC BY 4.0)}
\begin{document}
\maketitle

\noindent\emph{Source and licence.} A fractional Michael-Simon Sobolev inequality\\ on convex hypersurfaces. Version arXiv:2407.12439v3 (CC BY 4.0), distributed by arXiv under the Creative Commons Attribution 4.0 International licence, http://creativecommons.org/licenses/by/4.0/, as recorded in the arXiv licence metadata for this exact version and shown on the abstract page https://arxiv.org/abs/2407.12439v3. Full licence notice: \texttt{licenses/arxiv-2407.12439-CC-BY-4.0.txt}.\par\smallskip

\noindent\textit{Editorial scope.} Counted unit: the article's theorem theorem (source0.tex lines 304--312). The article's complete original proof of it is delivered with it (source0.tex lines 452--486, one \texttt{proof} environment of 35 lines, which the article places away from the statement). No mathematics is written, repaired or re-derived: the statement and the proof are the article's own lines, in the article's own notation, and no retained line is substituted or re-flowed.  Retained uncounted as the source-local context the proof needs: lines 370--385.  Nothing was omitted from any retained span, so the package carries no omission record.  No retained line contains a citation, so the wrapper carries no bibliography and no bibliography entry.  No retained line contains a figure environment, an image-inclusion command or drawing code, so no image is delivered and the package declares no asset dependency.  Editorial formatting only, in addition: an AMS \texttt{amsart} preamble that restates the article's own declarations and macros used by the retained lines, namely the environments \texttt{theorem}, \texttt{lemma} on separate counters, together with the standard packages those lines need.  The retained environments print in this document as Theorem 1, Lemma 1 in order of appearance, whereas the article prints the same items with the numbering of its own full document.  The complete native source of the article as distributed in the arXiv e-print for this exact version is bundled unchanged as \texttt{sources/162454/source0.tex}.
\begin{lemma}
 \label{gy:lemma:split into Mi}
  Let ~$\epsilon>0$ and ~$\Omega\subset\R^{n+1}$ be a convex open set. Then there exists  at most ~$2(n+1)$ closed sets ~$W_i\subset\R^n$ with nonempty interior, ~$i=1,\ldots,2(n+1)$, and Lipschitz functions ~$g_i:W_i\to \R$, such that the graphs of $g_i$
  $$
    \{(z,g_i(z)):\, z\in W_i\}=:M_i\,,
  $$
  satisfy that ~$\partial\Omega$ is the union 
  $$
    \partial\Omega=\bigcup_{i=1}^{2(n+1)}M_i\qquad\text{modulo rotations of each $M_i$} 
  $$
  up to a set of zero measure,
  and
  $$
    |\nabla g_i|\leq \sqrt{n+\epsilon}\quad\text{ in }W_i\,.
$$
\end{lemma}

\begin{theorem}
    \label{th:intro:key theorem}
     Let~$\Omega\subset\R^{n+1}$ be an open convex set and ~$\beta \in [0,n)$. Then, there exists a constant~$C$ depending only on~$n$ and $\beta$  such that for any measurable subset $E$ of $\delomega$,

     \begin{equation}
         \label{eq:intro: main estimate}
\int_{E}|x|^{-\beta}dx \leq C |E|^{\frac{n-\beta}{n}}.
     \end{equation}
\end{theorem}

\begin{proof}[Proof of Theorem \ref{th:intro:key theorem}]
Let ~$\epsilon$, ~$W_i$, ~$g_i$ and ~$M_i$ be as in Lemma \ref{gy:lemma:split into Mi}. We can chose $\epsilon=1$ for instance, but any $\epsilon>0$ does the job. For a measurable set ~$E\subset \partial\Omega$ we define ~$E_i:=E\cap M_i$. We shall also abbreviate the parametrizations ~$z\in W_i\mapsto(z,g_i(z))=\varphi_i(z)\in M_i$ and set ~$F_i=\varphi^{-1}(E_i)$. With these abbreviations and using Lemma ~\ref{gy:lemma:split into Mi} we have 
\begin{align*}
    \int_{E_i}\frac{1}{|x|^{\beta}}dx
    =&
    \int_{F_i}\frac{\sqrt{1+|\nabla g_i(z)|^2}}{(|z|^2+g_i^2(z))^{\beta/2}}dz
    \leq
    \sqrt{1+n+\epsilon}\int_{F_i}\frac{1}{|z|^{\beta}}dz
    \leq
    \sqrt{1+n+\epsilon}\int_{F_i^*}\frac{1}{|z|^{\beta}}dz,
\end{align*}
where ~$F_i^*$ is the Schwarz rearrangement of ~$F_i$ in $\R^n$, i.e., ~$F_i^*=\{z\in\R^n:\,|z|<R\}$ such that $\omega_n R^n=|F_i|$. Hence we obtain (recall the notation $\alpha_{n-1}=n\omega_n$)
\begin{align*}
        \int_{E_i}\frac{1}{|x|^{\beta}}dx
    \leq & \alpha_{n-1}\sqrt{1+n+\epsilon}\int_0^{R}r^{n-1-\beta}dr
    =\frac{\alpha_{n-1}\sqrt{1+n+\epsilon}}{n-\beta}\left(\frac{|F_i|}{\omega_n}\right)^{\frac{n-\beta}{n}}.
\end{align*}
Observe that ~$|F_i|\leq |E_i|$, which follows from integrating the inequality ~$1\leq \sqrt{1+|\nabla g_i|^2}$ over ~$F_i$, and thus we have proven that 
\begin{equation}
 \label{eq:gy:estimate for Ei}
   \int_{E_i}\frac{1}{|x|^{\beta}}dx\leq C_1 |E_i|^{\frac{n-\beta}{n}},\quad
   \text{ for }C_1=\frac{n\sqrt{1+n+\epsilon}}{n-\beta} \omega_n^{\frac{\beta}{n}}.
\end{equation}

Theorem \ref{th:intro:key theorem}  now follows by summing up the estimates of \eqref{eq:gy:estimate for Ei} over ~$i=1,\ldots, 2(n+1)$. We obtain from ~\eqref{eq:gy:estimate for Ei} and the generous estimate ~$|E_i|\leq |E|$ that 
\begin{align*}
    \int_{E}\frac{1}{|x|^{\beta}}dx\leq & \sum_{i=1}^{2(n+1)}\int_{E_i}\frac{1}{|x|^{\beta}}dx\leq C_1\sum_{i=1}^{2(n+1)}|E_i|^{\frac{n-\beta}{n}}
    \leq 
    2(n+1)C_1 |E|^{\frac{n-\beta}{n}},
\end{align*}
which proves the theorem with 
~$
  C=2(n+1)C_1
$.
\end{proof}

\end{document}
